\documentclass[11pt,a4paper]{article}
\usepackage[utf8]{inputenc}
\usepackage[T1]{fontenc}
\usepackage{amsmath,amssymb,amsthm}
\usepackage{graphicx}
\usepackage{booktabs}
\usepackage{hyperref}
\usepackage{geometry}
\geometry{margin=2.5cm}

\title{\textbf{Methodological Obstacles, API Deprecations, and Empirical Null-Results in Weak-Drive Coherence Anchoring on Superconducting Processors}}

\author{
  \textbf{M\'at\'e R\'obert}\thanks{Corresponding author: \texttt{silentnoisehun@gmail.com}}\\
  \textit{Independent Quantum Researcher}\\
  \textit{Budapest, Hungary}
}
\date{November 2026}

\begin{document}
\maketitle

\begin{abstract}
\noindent \textbf{Context:} The ``Quantum Anchor'' hypothesis postulated that continuous or periodic off-resonant weak microwave drives could dynamically anchor quantum phase and suppress environmental dissipation ($T_1$/$T_2$ decay) during idle times on transmon architectures. \\
\textbf{Methodological Hurdles:} Originally formulated for low-level analog pulse schedules (\texttt{qiskit.pulse}, \texttt{meas\_level=0}), experimental validation faced critical API deprecations across commercial cloud platforms, notably the phase-out of direct OpenPulse access by IBM and execution entitlement barriers on IQM platforms. Translating the concept into native circuit primitives revealed severe theoretical pitfalls, including amplitude-scaling oversights where nominal ``weak'' perturbations approached full $\pi$-pulse rotations. \\
\textbf{Empirical Results:} We conducted empirical evaluations on a 156-qubit IBM Heron processor (\texttt{ibm\_marrakesh}) and a 20-qubit IQM processor (\texttt{garnet}). On the Heron platform, single-qubit state survival reached $\psi(37\,\text{ns}) = 0.331662$, while both the unperturbed baseline and the active anchor condition yielded 0\% clear retention, failing to demonstrate dissipative mitigation. On IQM Garnet, multi-register entangling configurations achieved high per-plane Bell fidelity (95.4--96.6\%) and an 8-qubit GHZ state with 83.01\% global coherence (Wilson 95\% CI [80.59, 85.18]\%). However, an interferometric 5-phase parameter sweep showed global coherence confined to 0.20--1.37\%, situated at or below the hardware noise floor ($\sim 0.68\%$). \\
\textbf{Retractions and Null Findings:} We explicitly retract an earlier reported positive result of ``100\% state preservation'', demonstrating that it was an artifact of circuit-level self-inversion (a tautological identity operation) rather than genuine physical preservation. Goodness-of-fit tests reject both the theoretical oscillatory modulation and flat baseline at $p < 0.05$. \\
\textbf{Conclusions:} We present this study as a rigorous methodological null-result. The physical anchor drive hypothesis remains unsupported by empirical data. We outline clear recommendations for cloud-based quantum characterization, highlighting the critical distinctions between virtual phase rotations, circuit tautologies, and true physical drive dynamics.
\end{abstract}

\vspace{1em}
\noindent \textbf{Keywords:} Quantum characterization, Superconducting transmons, Null results, API deprecation, Dynamical decoupling, Methodological pitfalls.

\section{Introduction}

In quantum information processing on superconducting transmon circuits, coherence preservation against environmental dephasing ($T_2^*$) and energy relaxation ($T_1$) remains an essential challenge. The Quantum Anchor hypothesis was conceived as a heuristic state-stabilization strategy: it suggested that applying a calibrated, off-resonant continuous or periodic microwave tone during idling intervals could constrain transmon drift to a preferred manifold without inducing wave-function collapse.

Testing this protocol required access to raw, un-demodulated analog measurements and sub-gate microwave pulse generation. However, cloud-hosted quantum computing has rapidly standardized around rigid, gate-level compilation pipelines and abstracted execution primitives. In this article, we report the complete experimental trajectory of the Quantum Anchor protocol: from its original pulse-level formulation through multiple API deprecations and hardware trials, culminating in unambiguous empirical null results and the retraction of a circuit-level false positive.

The publication of well-documented negative results, failed protocols, and software deprecation impacts is crucial for avoiding duplicated community effort and unmasking systematic measurement artifacts in quantum empirical research.

\section{Evidence Classification Framework}

To ensure complete transparency and prevent over-interpretation of simulation models, all empirical and analytical claims throughout this investigation are categorized under a strict five-tier evidence hierarchy:

\begin{table}[htbp]
\centering
\begin{tabular}{lll}
\toprule
\textbf{Marker} & \textbf{Classification} & \textbf{Definition} \\
\midrule
$\checkmark$ & \textbf{PROVEN (Classical)} & Derived analytically or verified via offline classical simulation. \\
$\mathfrak{H}$ & \textbf{PROVEN (Hardware)} & Replicated on physical quantum processors with verifiable calibrations. \\
$\triangle$ & \textbf{UNVERIFIED} & Model assumption lacking empirical confirmation. \\
$\times$ & \textbf{FALSIFIED / NULL} & Hypothesized behavior rejected by physical hardware measurement. \\
$\mathbb{R}$ & \textbf{RETRACTED} & Prior positive finding shown to stem from systematic error or tautology. \\
\bottomrule
\end{tabular}
\caption{Evidence classification hierarchy applied across all experimental stages.}
\label{tab:evidence_grades}
\end{table}

\section{Methodological Barriers and API Deprecations (F1--F5)}

The transition from theoretical formulation to physical execution encountered five structural barriers:

\begin{description}
  \item[\textbf{F1 --- Removal of Qiskit Pulse:}] The original experimental design depended directly on \texttt{qiskit.pulse} scheduling (\texttt{pulse.drive\_channel}, \texttt{pulse.play}). IBM officially deprecated low-level OpenPulse control for production user-tier backends in 2024, removing direct pulse programming in early 2025. This rendered the original protocol inaccessible on public cloud QPUs.
  \item[\textbf{F2 --- Primitive Extraction Limits:}] In current Qiskit Runtime architectures (\texttt{SamplerV2} and \texttt{EstimatorV2}), the legacy parameters \texttt{meas\_level=0} (raw integrated IQ vector output) and \texttt{meas\_return='avg'} are non-existent. Output streams are restricted to post-discriminator digital bitstrings (\texttt{meas\_level=2}), obscuring analog intermediate states.
  \item[\textbf{F3 --- The Amplitude Scaling Oversight:}] Initial analytical specifications designated an anchor amplitude $\text{amp} = 0.08$ with a pulse duration $\tau = 37\,\text{ns}$ as a ``weak non-destructive perturbation''. However, integrating the effective Rabi rotation angle:
  \begin{equation}
    \theta_{\text{Bloch}} = \text{amp} \times \tau = 0.08 \times 37 = 2.96\,\text{rad} \approx 0.942\,\pi
  \end{equation}
  A $2.96\,\text{rad}$ rotation represents $94.2\%$ of a complete state inversion ($\pi$-pulse). Consequently, the drive was not perturbative; it acted as a strong, near-complete spin inversion.
  \item[\textbf{F4 --- Unphysical Zero-Dissipation Limit ($\gamma=0$):}] Theoretical benchmark models posited an idealized dissipation-free condition ($\gamma=0$). Because all physical superconducting transmon qubits operate under finite dielectric loss and Purcell decay ($T_1 \sim 80\text{--}150\,\mu\text{s}$), this baseline cannot be realized on physical hardware.
  \item[\textbf{F5 --- Cloud Entitlement Barriers:}] On the IQM Resonance platform (IQM Garnet 20Q), low-level pulse sweeps require enterprise-level permissions. Submissions targeting the pulse-level \texttt{SweepDefinition} endpoint were rejected with authorization denial errors (\texttt{Personal account does not have pulse-level access enabled}), restricting evaluations strictly to circuit-level operations.
\end{description}

\section{IBM Hardware Evaluation: ibm\_marrakesh (156Q Heron r2)}

\subsection{Protocol Migration to Fractional Gates}
To circumvent the deprecation of \texttt{qiskit.pulse}, we migrated the protocol into fractional gate rotations ($\text{RX}(\theta), \text{RZ}(\phi)$) combined with open-system Lindblad dynamics simulation via \texttt{qiskit-dynamics}.

Measurements were executed on the 156-qubit IBM Heron r2 processor (\texttt{ibm\_marrakesh}) with 2000 shots per circuit setting. Single-qubit state survival was tracked under varying parameter depths ($D_0$ through $D_8$).

\begin{table}[htbp]
\centering
\begin{tabular}{lccc}
\toprule
\textbf{Configuration / Condition} & \textbf{Observed Metric} & \textbf{Value} & \textbf{Status} \\
\midrule
$D_0$ initial layer & Mean fidelity & $97.40\%$ & $\mathfrak{H}$ \\
$D_4$ intermediate depth & Mean fidelity & $93.10\%$ & $\mathfrak{H}$ \\
$D_8$ deep fractional layer & Mean fidelity & $89.40\%$ & $\mathfrak{H}$ \\
Fidelity Preservation criterion ($D_0 - D_8 < 2\%$) & Evaluated boolean & \textbf{False} & $\times$ \\
Single-qubit state survival $\psi(37\,\text{ns})$ & Fitted amplitude & $0.331662$ & $\mathfrak{H}$ \\
Unperturbed baseline ($\gamma=0$ control) & Clear state retention & \textbf{0.00\%} & $\times$ \\
Active weak-drive condition ($\gamma=0.5$) & Clear state retention & \textbf{0.00\%} & $\times$ \\
\bottomrule
\end{tabular}
\caption{Empirical results on \texttt{ibm\_marrakesh} showing monotonic gate-depth decay and identical failure in state retention across control and active drive conditions.}
\label{tab:ibm_results}
\end{table}

\subsection{Retraction of False Positive: The Circuit-Tautology Trap}
An earlier preliminary iteration of our code reported ``$100\%$ clear state preservation'' under the $\gamma=0$ model, which was initially interpreted as confirmation of anchor stabilization.

Subsequent static circuit decomposition revealed this to be a \textbf{circuit tautology}:
The $\gamma=0$ implementation applied forward rotations followed immediately by their exact algebraic inverse ($\hat{U}^{\dagger}\hat{U} = \hat{I}$). The apparent perfect fidelity was simply an identity operation. When non-inverting dissipation operators were injected into both arms, both the control condition ($\gamma=0$) and active drive condition ($\gamma=0.5$) decayed to \textbf{0\% clear retention} (Table~\ref{tab:ibm_results}). We formally retracted the claim of state preservation.

\section{IQM Hardware Evaluation: Garnet 20-Qubit Processor}

\subsection{Experimental Architecture}
To assess whether multi-register state correlations could stabilize phase drift without analog pulse access, an 8-qubit register was partitioned into four 2-qubit partitions ($q_0\text{--}q_1, q_2\text{--}q_3, q_4\text{--}q_5, q_6\text{--}q_7$). Native gates on IQM Garnet include $\{\text{CZ}, \text{PRX}, \text{RZ}\}$.

A virtual Z-rotation $\text{RZ}(\theta)$ was introduced with $\theta = 0.4398\,\text{rad}$ ($2\pi \times 4.11\,\text{GHz} \times 37\,\text{ns} \pmod{2\pi}$). Crucially, this operation is realized entirely through phase-offset tracking in control software; it deposits zero physical energy onto the resonator drive line.

\subsection{Multi-Register GHZ Entanglement (V2 Topology)}
An inter-partition entangling spine connecting $q_0 \leftrightarrow q_2 \leftrightarrow q_4 \leftrightarrow q_6$ followed by intra-plane Bell extensions was measured on IQM Garnet (Job ID \texttt{01a1162c-717c-77e7-91d9-90ed16c0e591}):
\begin{itemize}
  \item Per-partition Bell fidelities demonstrated uniform performance across pairs: $P_0 = 95.41\%$, $P_1 = 96.09\%$, $P_2 = 96.19\%$, $P_3 = 96.58\%$.
  \item An 8-qubit entangled GHZ state reached \textbf{83.01\% global coherence} (850/1024 shots), with a Wilson 95\% confidence interval of $[80.59\%, 85.18\%]$.
\end{itemize}
This verified that clean multi-qubit coherent states can be established across the device. However, this benchmark reflected standard gate-synthesis performance rather than evidence of autonomous drive-induced stabilization.

\subsection{Interferometric Phase Sweep: Noise Dominance (V3 Protocol)}
Because projective Z-basis measurements on GHZ states are insensitive to relative phase shifts $\phi$, an interferometric Hadamard transform ($H^{\otimes 8}$) was applied immediately prior to readout across a 5-point phase parameter sweep ($\phi \in \{0.00, 1.26, 2.51, 3.77, 5.03\}\,\text{rad}$).

\begin{table}[htbp]
\centering
\begin{tabular}{ccccc}
\toprule
\textbf{Phase Index} & \textbf{Phase $\phi$ (rad)} & \textbf{Theoretical Signal} & \textbf{Measured Coherence} & \textbf{Wilson 95\% CI} \\
\midrule
0 & 0.0000 & 3.125\% & 1.27\% & [0.74, 2.16]\% \\
1 & 1.2566 & 1.078\% & 0.29\% & [0.10, 0.86]\% \\
2 & 2.5133 & 0.375\% & 0.20\% & [0.05, 0.71]\% \\
3 & 3.7699 & 0.375\% & 1.37\% & [0.82, 2.28]\% \\
4 & 5.0265 & 1.078\% & 0.29\% & [0.10, 0.86]\% \\
\bottomrule
\end{tabular}
\caption{Hardware measurement of 5-phase interferometric readout on IQM Garnet.}
\label{tab:iqm_sweep}
\end{table}

\noindent \textbf{Statistical Evaluation:}
\begin{enumerate}
  \item Global state survival measured between $0.20\%$ and $1.37\%$.
  \item The physical noise floor of the readout ($\sim 0.68\%$) equals or exceeds the predicted interferometric contrast.
  \item Goodness-of-fit testing against the theoretical $\cos(8\phi)$ model yields $\chi^2 = 30.72$ ($\text{df}=4, p = 6.4 \times 10^{-6}$), soundly \textbf{rejecting} agreement with theoretical predictions.
  \item Testing against an unmodulated flat noise model similarly yields $\chi^2 = 20.29$ ($\text{df}=4, p = 5.2 \times 10^{-4}$), indicating that fluctuations across shots reflect drift rather than coherent phase modulation.
  \item The empirical ordering diverged completely from theory: Phase 3 (theoretically $0.375\%$) registered the highest count ($1.37\%$), while Phase 0 (theoretically $3.125\%$) registered only $1.27\%$.
\end{enumerate}

\section{Discussion and Methodological Lessons}

The empirical findings from this investigation provide several critical takeaways for experimental quantum computing:

\begin{enumerate}
  \item \textbf{The Software-Hardware Disconnect in Cloud QPU Access:} The withdrawal of pulse-level APIs by major cloud providers fundamentally alters what can be experimentally evaluated. Attempts to emulate continuous physical stabilization using software-compiled circuit gates or virtual Z rotations (\texttt{RZ}) cannot capture analog non-linear drive effects.
  \item \textbf{The Circuit-Tautology Hazard:} In multi-stage quantum algorithms, especially those involving forward-backward error cancellation or dynamical recovery, sanity audits must prove that state survival is not simply the trivial consequence of unitary self-inversion ($\hat{U}^{\dagger}\hat{U} = \hat{I}$).
  \item \textbf{Statistical Rigor and Noise Floor Realism:} When theoretical interference signals drop below a few percent (e.g., $1/128 \approx 0.78\%$), standard shot allocations (1024 shots) cannot resolve true coherent oscillation from background state-preparation and measurement (SPAM) error.
\end{enumerate}

\section{Conclusion}

The Quantum Anchor protocol does not produce measurable dissipative mitigation or autonomous phase stabilization on physical transmon processors under gate-level access. By cataloging the structural blockers, API barriers, and empirical null results, we establish a definitive experimental baseline that clarifies the operational boundaries between theoretical open-system drive models and current cloud quantum infrastructure.

\section*{Data and Code Availability}
All raw bitstring records, calibration logs, execution job IDs, and offline verification scripts are openly available in the research repository: \url{https://github.com/silentnoisehun/quantum-anchor}.

\bibliographystyle{plain}
\begin{thebibliography}{99}
\bibitem{jote_aims}
van Dongen, N. N., \& Sikorski, J. (2021). The Journal of Trial and Error: A platform for publishing negative and inconclusive findings. \textit{Journal of Trial and Error}, 1(1), 1--4.

\bibitem{ibm_runtime}
IBM Quantum. (2024). \textit{Qiskit Runtime Primitives and Migration Guide}. IBM Documentation.

\bibitem{transmon_coherence}
Krantz, P., Kjaergaard, M., Yan, F., Orlando, T. P., Gustavsson, S., \& Oliver, W. D. (2019). A quantum engineer's guide to superconducting qubits. \textit{Applied Physics Reviews}, 6(2), 021318.

\bibitem{dynamical_decoupling}
Viola, L., Knill, E., \& Lloyd, S. (1999). Dynamical decoupling of open quantum systems. \textit{Physical Review Letters}, 82(12), 2417.
\end{thebibliography}

\end{document}
